\chapter{状态动力学的基本原理：目标、结构与受限更新}
\label{chp:first_principles}

目的交互主义确定了我们研究的对象：具有历史、目标和边界的系统，在环境反馈中持续生成可修订状态。本章将这一立场推进为可计算的动力学模型。我们所说的基本原理，不是从某个物理名称推出全部智能行为，而是建立推导能够成立的起点：状态必须带有类型与身份，评价必须说明目标及尺度，更新必须遵守可行域，结构改变必须留下事件记录，实现必须接受资源与时延限制。在此基础上，我们分别构造静态目标、连续更新、离散迭代和轨迹规划，说明哪些条件支持下降、哪些条件支持收敛，以及为什么这些数学性质仍不能替代现实任务成功。表示扩展与资源压缩、历史连续与目标转向、局部观测与全局控制等问题都在本章保留，但其机制改由明确的算子、接口和预算表达。HSF-HD 因此提供智能一般状态动力学的模型族；AgentOS 在下一理论层实现其中一部分模型，而不构成所有智能的唯一结构。

\section{双目视角：任务分辨率与资源预算的共同约束}
\label{sec:section_3791}

智能系统既要区分对任务重要的状态，又不能无限保存和计算所有差异。前一要求推动表示扩展，后一要求促使选择、压缩、延迟处理或卸载。我们把二者称为任务分辨率与资源预算的双目视角。这里不存在预先保证平衡的自然法则：表示过少可能漏掉风险，表示过多也可能导致过载、错误绑定和无效搜索。有效组织来自在具体任务、误差容限和资源条件下选择合适的状态变量，而不是把外部世界全部复制到内部。环境仍有未观测部分，内部也有只用于推演的候选，两者不能被“镜像”一词直接等同。

设环境对象为 $w\in\mathcal W$，观测为 $o\in\mathcal O$，编码器 $e$ 把观测及历史映射为表示 $r\in\mathcal R$，读出器 $d$ 给出任务判断。我们在声明的数据分布 $\mathcal D$ 上定义任务失真
\begin{equation}
 D(e,d)=\mathbb E_{\mathcal D}
 \bigl[\ell(d(e(o)),y)\bigr],
\end{equation}
其中 $y$ 是任务判据，$\ell$ 是非负损失。不同任务有不同损失：物体身份误配、路线时延、事实错误和违反权限不能未经解释都算成同一种误差。表示成本向量 $c(e,d)$ 分别记录内存字节、计算次数、时延和外部调用费用，预算向量为 $B$。设计问题是寻找满足 $D\le\varepsilon$ 且 $c\preceq B$ 的实现，或者在合法域内比较这些指标的帕累托前沿；若不存在这样的实现，我们必须报告误差要求与预算不相容。

局部结构网络负责对象身份、关系类型与依赖，全局分布表示负责语境、相似性和难以枚举的统计规律。表示扩展可以增加节点与边，也可以提高属性精度、扩大候选覆盖或延长历史范围，不能只用维度增长描述。例如机器人辨认“可承重梁”时，视觉向量帮助识别外形，局部结构保存材料批次、安装位置和检验记录。增加像素分辨率可能改善外形识别，却无法替代材料认证；增加关系节点可能完善施工依赖，也不能自动获得新的传感证据。每种扩展都应指出它减少哪一类任务失真，以及是否引入新的身份或来源风险。

压缩同样需要保留任务相关条件。若两个环境状态在当前任务下要求不同动作，却被编码成同一表示，读出器就无法仅凭该表示同时作出正确选择。因此压缩应保留动作相关的可区分性，而不是只追求重构平均误差小。杯子是否有裂纹可能仅占少量像素，却决定能否盛热水；文档中的否定词可能极短，却改变整条承诺。平均失真掩盖稀有高风险差异时，我们应增加分组指标、最坏情形检查或硬约束，避免把压缩效率误称为语义充分。

信息熵在这一问题中只描述指定随机变量的分布不确定性。它不等于内存占用、理解深度、活动强度或焦耳能耗。相同熵的编码可以具有不同检索延迟，相同存储量也可能保留不同任务价值。实际器件能耗当然限制实现，但从某个热力学下界不能直接推出应该遗忘哪条承诺或保留哪个属性。选择规则仍需任务损失、预算与授权共同决定。我们使用耗散理论理解开放实现需要资源交换，使用认知科学理解有限工作界面，却不把二者混成一种万能代价。

一个具体反例说明“越丰富越智能”并不成立。配送代理同时接入天气、交通、新闻和所有用户历史，若有限工作区缺乏身份绑定与优先规则，噪声可能挤出药品温控期限，导致更高失误率。相反，只保存地址和路线也可能遗漏接收授权。合理方案是维持当前任务的最小可靠记录，把低频资料置于外部记忆，并按需要召回。这不是永远选择最短描述，而是在声明的任务族上检验表示充分性，保留异常触发的扩展通道。

我们由此把形与质的联合表示落实到有类型的计算组织。节点属性可以作为离散纤维式的归属结构，但离散关系网络并不因此自动成为光滑流形。只有额外给出坐标图、转移映射和连续性条件，才可使用微分几何；只有给出概率模型与可辨识参数，才可使用信息几何。原有膜面意象帮助读者看见“约束塑造可行区域”的直觉，实际推导则依赖误差、接口与预算。

\begin{figure}[htbp]
 \centering
 \includegraphics[width=0.70\linewidth,height=0.40\textheight,keepaspectratio]{image/geom_physic.png}
 \caption{受限状态空间的辅助意象。起伏与路径示意目标、结构和边界的共同约束；画面不代表测得的几何曲率或真实物理力。}
\end{figure}

\section{目标泛函：由观测、结构与目标构造更新方向}
\label{sec:section_8054}

为了讨论状态怎样变化，我们先建立一个可完整演算的有限维特例。固定结构区间内，令 $x\in\mathbb R^n$ 为经过参考尺度归一化的连续状态，$y\in\mathbb R^m$ 为归一化观测，$C\in\mathbb R^{m\times n}$ 为观测读出矩阵。矩阵 $W\in\mathbb R^{m\times m}$ 对称半正定，用于表达观测残差的权重；它不是在未校准条件下自动等于置信度。关系约束算子为 $A\in\mathbb R^{p\times n}$，其第 $j$ 行描述第 $j$ 个可解释的关系残差。目标状态为 $r\in\mathbb R^n$，目标权重 $Q$ 对称半正定。各对象都固定在同一身份和版本上，否则矩阵乘法虽合法，语义仍可能错配。

我们定义无量纲目标
\begin{equation}
 V(x)=\frac12(Cx-y)^\top W(Cx-y)
 +\frac{\lambda}{2}\|Ax\|^2
 +\frac12(x-r)^\top Q(x-r),\qquad \lambda\ge0.
\end{equation}
第一项衡量当前表示与观测的失配，第二项惩罚指定关系的不一致，第三项表达可允许权衡的目标偏差。这是模型定义，不能由“目标驱动”四个字唯一推出。权限、安全与身份规则应另外进入可行域 $\mathcal K$，不能被很小的惩罚权重抵消。若观测冲突，最小化第一项得到的是某种折中估计，不是自动发现唯一事实；若目标不可达，第三项也只记录差距，不授权违反边界。

取任意扰动 $v\in\mathbb R^n$，把 $x+\epsilon v$ 代入并展开一次项，得到
\begin{equation}
 \left.\frac{d}{d\epsilon}V(x+\epsilon v)\right|_{\epsilon=0}
 =v^\top\!\left[C^\top W(Cx-y)+\lambda A^\top Ax+Q(x-r)\right].
\end{equation}
因为该式对所有 $v$ 成立，欧氏梯度就是方括号中的向量，Hessian 为
\begin{equation}
 H=C^\top WC+\lambda A^\top A+Q.
\end{equation}
对任意 $v$，有 $v^\top Hv=(Cv)^\top W(Cv)+\lambda\|Av\|^2+v^\top Qv\ge0$，所以这个特例是凸二次模型。当且仅当非零方向不会同时落入三个惩罚的零空间时，$H$ 正定，才能在无约束域得到唯一极小点。特别地，关系图连通不足以保证唯一性：仅有差分约束时常量方向仍可能不可辨识。

展开 $V$ 可写为 $\frac12x^\top Hx-b^\top x+\mathrm{const}$，其中 $b=C^\top Wy+Qr$。若 $H$ 正定，完成平方或解线性方程得到 $x^*=H^{-1}b$；若 $H$ 奇异，先检查 $b$ 是否位于 $H$ 的像空间，再讨论非唯一解。我们不以添加任意微小正则就宣称问题已经认识论上可辨识；正则可以选择一个解，却可能把未经证据支持的先验写入结果。报告时应注明哪些方向由观测确定，哪些方向主要由结构或目标约束选择。

这里的结构项也不能等同于逻辑真理。若 $A$ 只是无向图的关联矩阵，$\|Ax\|^2$ 只要求相邻属性平滑，不保证因果判断、类型推理或事实正确。互相矛盾的证词不应为了降低差分而被强行平均；不同角色的数值也不能只因有边连接就混合。对于关系类型不同的对象，我们需要先给出转换算子及残差含义，再构造每行 $A$。强连贯的错误叙述同样可能拥有很小结构代价，因此观测、结构和目标必须分别保留诊断指标。

目标项表达“未来约束”时，未来不会反向施加真实力。它来自当前时刻可用的目标记录、预测模型和授权信息。预测器可能误估未来结果，目标也可能改变；更新方向只是当前模型对候选状态的排序。实际策略仍需通过动作接口影响环境，再用后果验证。把这个梯度称为引力可以帮助直观想象方向，却无法增加其经验有效性，更不能证明所有智能遵循同一标量目标。

这一构造还展示了模型族的开放性。非线性观测可把 $Cx$ 改成可微读出，离散结构编辑可把矩阵版本改成事件，多个目标可用偏序或约束族代替单个 $Q$。每次推广都要重新证明存在性、正定性和适用条件，不能沿用线性特例的结论。HSF-HD 的基本原理因此是一套构造和检查更新模型的方法；明确给出的二次目标是可复现示例，其余智能实现需要各自的任务定义与证据。

\section{连续更新：完整导数、历史延续与扰动边界}
\label{sec:state_continuous_update}

目标函数给出局部比较方向，更新规则才决定实际轨迹。令时间 $t$ 以声明的采样周期为单位归一化，状态仍为无量纲向量，取对称正定的更新度量 $M\in\mathbb R^{n\times n}$。我们选择一阶更新模型 $M\dot x=-\nabla_xV$。$M$ 控制不同方向的更新尺度，可以来自预条件、参数敏感度或人工设计，但不因此是思维质量。固定 $V$ 与 $M$ 时，由链式法则得到 $\dot V=\nabla V^\top\dot x=-\nabla V^\top M^{-1}\nabla V\le0$。这个数学结果说明模型目标不增加，尚未说明状态有界、梯度趋零或任务成功；这些结论各需补充条件。

对前节正定二次目标，令误差 $e=x-x^*$。更新满足 $M\dot e=-He$，并可用 $\frac12e^\top Me$ 衡量误差。其导数为 $-e^\top He$；由于 $H$ 正定且 $M$ 正定，可以用最小与最大特征值比较二次型，得到误差指数衰减。若 $H$ 仅半正定，则未受约束方向未必衰减，更不能从目标下降推出唯一状态。若目标非凸，则下降轨迹可能进入局部驻点；若可行域不闭或目标不下有界，也不能借这一式保证过程终止。

现实系统的目标、结构、输入与度量会变化。把前节各项写成时变对象，并记 $z=C_tx-y_t$、$a=A_tx$、$e=x-r_t$，完整全导数为
\begin{align}
 \frac{dV_t}{dt}
 ={}&\nabla_xV_t^\top\dot x
 +z^\top W_t(\dot C_tx-\dot y_t)
 +\frac12 z^\top\dot W_tz \notag\\
 &+\frac{\dot\lambda_t}{2}\|a\|^2
 +\lambda_ta^\top\dot A_tx
 -e^\top Q_t\dot r_t+\frac12e^\top\dot Q_te.
\end{align}
这里把自变量 $x$ 的变化放在首项，其余全部属于模型或目标的显式变化。如果采用时变度量，它通过更新规则改变 $\dot x$；若还把 $M_t$ 写进误差评价或速度评价，其导数也必须出现。矩阵来自离散图编辑时不应强行求导，而要在事件前后分别计算目标，并记录目标可能跳变。省略这些项会把环境变化造成的目标增加误判为控制器失败。

考虑扰动 $d(t)$ 下的模型 $M_t\dot x=-g+d$，其中 $g=\nabla V_t$。记所有显式变化项之和为 $\partial_tV_t$，则
\begin{equation}
 \dot V_t=-g^\top M_t^{-1}g+g^\top M_t^{-1}d+\partial_tV_t.
\end{equation}
在 $M_t$ 对称正定时，令 $\|v\|_{M_t^{-1}}^2=v^\top M_t^{-1}v$，由平方非负或柯西不等式得到
\begin{equation}
 \dot V_t\le-\frac12\|g\|_{M_t^{-1}}^2
 +\frac12\|d\|_{M_t^{-1}}^2+\partial_tV_t.
\end{equation}
因此只有下降项足以覆盖扰动和目标漂移时，才能宣称这一时刻下降。目标变动越快，系统越可能只能跟踪一个邻域；该邻域大小还依赖强凸性、扰动界与变化速度，不能无条件声称必然收敛。

历史延续可以另用二阶更新表达。固定对称正定 $M$、半正定 $D$ 时，设 $M\ddot x+D\dot x+\nabla V(x)=d$，速度 $\dot x$ 是增加的算法状态。定义评价 $\mathcal E=V(x)+\frac12\dot x^\top M\dot x$，直接求导并代入方程得到 $\dot{\mathcal E}=-\dot x^\top D\dot x+\dot x^\top d$。若 $M,V$ 时变，还需加 $\frac12\dot x^\top\dot M\dot x+\partial_tV$。这一式区分了阻尼、输入与历史速度；结构梯度并不自动提供耗散，零阻尼模型可能持续振荡。

二阶模型的速度记忆可以解释更新不立即转向的算法现象，但不能据此把悲伤、习惯或胆量换算成单一质量常数。人的心理过程需要额外观测和经验检验，软件动量项则需要测试其超调、约束违反与延迟反馈风险。例如目标突然从“快速搬运”改成“紧急停止”，沿旧方向延续的速度可能降低平滑性代价，却违反安全边界；正确实现必须允许硬约束触发中止。历史连续性是受条件控制的设计目标，不是不可违背过去的普遍命题。

连续描述的价值在于提供局部分析工具。它帮助我们识别更新尺度、扰动来源和目标漂移，给出可检验的不等式，而不是把全部认知变化连续化。读出、承诺提交、结构删除和权限变更仍然是离散事件。我们在固定身份与结构区间分析轨迹，再用事件接口连接区间，才能同时保留数学严谨性与实际系统的混合特征。

可行域也必须参与连续模型。无约束梯度方向可能越过容量、权限或安全边界，即使它使目标下降。对于固定闭凸可行集，我们可以将方向限制在局部切锥并检查边界保持；对于离散授权与身份规则，则需要事件守卫和合法提交，而不能用一个连续投影掩盖。若可行域自身收缩，原来合法的状态也可能变得不可行，这时应触发退避、停止或重新协商。所有连续下降结论都应附带轨迹仍在模型有效域内这一条件，域外行为由边界机制处理。

\section{离散实现与轨迹变分：两类问题的推导条件}
\label{sec:state_discrete_variation}

数字系统最终通过离散步骤运行，连续方程的下降性质必须传递到具体积分器。对固定、可微且梯度为 $L$-Lipschitz 的目标 $V$，取更新 $x^+=x-\eta g$，其中 $g=\nabla V(x)$、$\eta>0$。沿更新线段积分，有 $V(x^+)-V(x)=\int_0^1\nabla V(x-\tau\eta g)^\top(-\eta g)\,d\tau$。把初始梯度分离，利用 Lipschitz 条件界定剩余项，得到
\begin{equation}
 V(x^+)\le V(x)-\eta\left(1-\frac{L\eta}{2}\right)\|g\|^2.
\end{equation}
因此 $0<\eta<2/L$ 且梯度非零时严格下降。条件必须在线段经过的区域成立；这不是任意步长的保证，也不保证非凸目标的全局最优。若目标在更新后改变，还要单独加上 $V_{k+1}(x^+)-V_k(x^+)$，不能继续声称跨版本单调。

若状态受非空闭凸集合 $\mathcal K$ 约束，采用欧氏投影 $x^+=P_{\mathcal K}(x-\eta g)$。假设初值 $x\in\mathcal K$，由投影的最优性条件可得 $g^\top(x^+-x)\le-\|x^+-x\|^2/\eta$。与相同的光滑上界结合，
\begin{equation}
 V(x^+)\le V(x)-\left(\frac1\eta-\frac L2\right)\|x^+-x\|^2.
\end{equation}
这为约束内更新提供下降条件。若集合由离散权限、图类型和组合规则形成，通常不是凸集，欧氏投影可能多值或计算困难；我们应使用专门的合法编辑、搜索或拒绝机制。把符号约束画成平滑边界不能消除这种差异。

对于二次目标，误差迭代为 $e_{k+1}=(I-\eta H)e_k$。若 $H$ 对称正定，在其特征基上每个方向乘以 $1-\eta\lambda_i$，所以 $0<\eta<2/\lambda_{\max}(H)$ 使所有方向模小于一并线性收缩。异步读写、延迟梯度或近似算子会改变该矩阵，必须重新分析。工程日志应记录目标版本、步长、梯度来源与拒绝原因，避免把采用旧目标的更新同最新目标混在一条“下降曲线”里。

轨迹规划是另一类问题。它选择整段路径，而不是由局部更新方程自动生成路径。给定固定端点 $x(0)=x_0$、$x(T)=x_T$，我们定义平滑轨迹代价
\begin{equation}
 \mathcal J[x]=\int_0^T
 \left[V_t(x(t))+\frac12\dot x(t)^\top R_t\dot x(t)\right]dt,
\end{equation}
其中时间、状态和积分尺度已归一化，$R_t$ 对称正定且可微。速度项惩罚改变过快的规划，不必解释为物理动能。对端点为零的扰动 $v(t)$，一次变分为 $\delta\mathcal J=\int_0^T[\nabla V_t^\top v+\dot x^\top R_t\dot v]dt$。分部积分消去边界项，得到
\begin{equation}
 \delta\mathcal J=\int_0^T
 \left[\nabla V_t-\frac{d}{dt}(R_t\dot x)\right]^\top v\,dt.
\end{equation}
若驻值对所有允许扰动成立，必要条件便是 $\frac{d}{dt}(R_t\dot x)=\nabla V_t$，也即 $R_t\ddot x+\dot R_t\dot x=\nabla V_t$。符号与前节下降动力学不同，原因是一个求整段轨迹的驻值，一个规定在线因果更新；二者不能凭方程相似而混用。

驻值条件本身还不是极小值证明。二次变分为 $\int_0^T[v^\top\nabla^2V_tv+\dot v^\top R_t\dot v]dt$，在 $V_t$ 凸且 $R_t$ 一致正定时为非负，并可用于证明已有驻值路径的最小性质；路径存在还要给出正则性、下有界性与合适函数空间条件。若终点自由，分部积分的末端项不消失，需要相应的自然边界或终端代价条件。若外部动作满足控制方程，还应把控制变量及动力学约束加入规划，不能假定任意内部路径都可现实执行。

以机器人操作工具为例，规划器可以利用速度惩罚减少不必要改向，但真实设备还有扭矩、接触、延迟和障碍限制。优化路径可能在内部编码空间平滑，却在外部碰撞；这时需要通过动作读出和模拟重新验证。软件代理也类似，连续改变一个内部意图向量不能代替数据库事务或代码编译。优化算法的收敛、内部轨迹的连贯和执行后果的成功是三个不同证据层次。

我们因此保留变分方法的价值，同时限定其角色。它可以导出声明了端点、目标与速度代价的轨迹规划条件，也可以帮助检查系数、符号与边界项；它不证明所有智能遵循最小作用量，不证明当前目的能从未来施力，也不证明驻值轨迹必然是测地线。对于一般 HSF-HD 模型，在线更新、离散编辑和离线规划可以共同存在，但每个模块都需要明确输入、状态和验证契约。

\section{多尺度控制：局部观测、运行状态与目标协调}
\label{sec:section_5140}

观测接口、运行更新与长期目标经常具有不同的频率和作用范围，因此分层控制是一种重要工程选择。局部环处理传感、输入清洗和快速反馈，中层环维护工作状态、候选和任务连续性，高层环修订目标、承诺、预算和慢变结构。分层不意味着各层正交，也不能由资源有限直接证明所有智能必须具有同一种层数。一个简单系统可以用单层控制满足任务，一个复杂系统也可能采用事件驱动或分布式协商。我们讨论的是何时分层能够减少延迟与重复计算，以及它需要怎样的接口才能避免失控。

令快状态 $x(t)\in\mathbb R^n$，慢参数 $z(t)\in\mathbb R^q$，在归一化时间上建立模型
\begin{equation}
 \dot x=f(x,z,o),\qquad \dot z=\epsilon g(x,z,q),
 \qquad 0<\epsilon\ll1.
\end{equation}
小参数只表达相对更新速度。若对每个固定 $z$，快系统存在唯一稳定平衡 $x^*(z)$，并具有一致的收敛速率，若 $x^*(z)$ 随 $z$ 足够平滑，慢更新和观测变化又有界，我们才有理由用 $g(x^*(z),z,q)$ 近似慢动力学。缺少唯一性、快环稳定性或尺度分离时，慢环可能收到不同吸引域的状态，近似就会失败。所谓低频控制更持久，来自它修改的约束会被后续多步重复使用，而不是频率本身提供特殊力量。

我们可以用一个局部误差模型看见快慢耦合的影响。假设 $x^*(z)$ 可微，设 $e=x-x^*(z)$，则 $\dot e=f(x,z,o)-D_zx^*(z)\dot z$。若快环围绕平衡的收缩率为 $a>0$，并能界定观测变化及慢环驱动的幅度为 $b$，误差范数可在相应有效域满足 $\frac{d}{dt}\|e\|\le-a\|e\|+b$，从而通过积分因子得到 $\|e(t)\|\le e^{-at}\|e(0)\|+(b/a)(1-e^{-at})$。这个界说明快速收缩与缓慢目标变化共同决定跟踪误差，不能只检查其中一层稳定便宣布整体稳定。

局部观测既不能被视为绝对钉扎的真值，也不能任由高层目标改写。输入应保存来源、时间、校准状态和有效期；系统根据误差模型形成估计或置信集合。可靠身份与权限记录可以构成硬约束，带噪测量则通常形成软残差或区间约束。若估计与目标冲突，系统要检查观测失效、目标不可行、模型遗漏和身份错配，而不是把冲突统一叫作“惊奇”。统计惊奇若用负对数概率定义，需要明确概率模型；它与欧氏残差或实际设备应力不是同一个量。

在图上的局部观测例子更清楚。令 $L$ 为固定无向非负权图的拉普拉斯，状态分为观测边界 $x_B$ 与待估内部 $x_I$。最小化 $\frac12x^\top Lx$ 且固定 $x_B$ 时，对内部求导得到 $L_{II}x_I=-L_{IB}x_B$。只有每个相关连通分量都与边界相连等条件保证 $L_{II}$ 正定时，才得到唯一调和延拓。这个计算可以填补指定平滑属性，却不能从温度或嵌入平滑推导语义理解。目标项或新观测加入后方程改变，未知边界也不能被固定观测假设掩盖。

上下层接口要传摘要与约束。快环向任务环报告对象身份、异常、置信范围、资源余量与有效时间，而不是上传未经筛选的全部原始信号；任务环向慢环报告跨步骤偏差、未履行承诺与修改候选。慢环下传预算、阈值和允许策略，不应以一个模糊全局信号覆盖局部安全中止。摘要压缩若漏掉稀有风险或破坏时间顺序，就会让高层在错误状态上优化。接口质量必须通过回放、故障注入和跨版本检查评估。

AgentOS 可把这类分层具体化为 $h(t)$、$E$ 与 $\Theta$ 的多尺度更新，以及 SPC 的自我相关状态摘要、TCE 的元控制、CCM 的复杂度监测、Boundary 的交换限制和 Offloading 的外部重构。全息自我 $\Psi$ 在这里保留为该架构的自我相关组织变量，不再用作任意复数认知场的别名。结构网络与全局分布状态仍通过身份绑定联合；我们不从前额叶、海马体或膜结构的相似叙述推出计算机制相同。

两个原图中的曲面、核心和边界可以帮助组织这一层级意象。图上带有历史物理名称的文本应结合这里的计算解释阅读：边界对应观测与执行接口，介质对应实现与状态组织，高层对应较慢目标协调，自我对应工程实例的自我相关状态。图像不是层级必然性的证明；模型的依据仍是尺度、作用域、误差及接口条件。

\begin{figure}[htbp]
 \centering
 \includegraphics[width=0.98\linewidth,height=0.46\textheight,keepaspectratio]{image/tpmodel.png}
 \caption{环境接口、运行组织与目标协调的原有意象图。图中曲面及物理名称仅辅助阅读；计算定义分别由观测、状态更新和多尺度控制接口给出。}
\end{figure}

\section{实现约束：传播时延、保留机制与探索尺度}
\label{sec:implementation_computational_constraints}

计算模型必须落在具体实现上，但其限制不能归结为三个万能介质常数。网络传播速度、排队时延、调用配额、内存容量、存储可靠性和实际功耗共同影响可用状态更新。不同实现即使执行同一数学映射，也可能有不同故障分布和预算；同一设备在负载变化后也会呈现不同延迟。我们把实现约束写成需要测量的接口契约，分别说明量纲、测量窗口和适用工作负载，不把“认知光速”当作语义图距离与秒之间无需定义的转换常数。

对于由多个阶段组成的观测到动作链，令 $\tau_i$ 为第 $i$ 个阶段的处理延迟，$\delta_i$ 为通信延迟，$q_i$ 为排队等待。串行链至少需要 $\sum_i(\tau_i+\delta_i+q_i)$；有并行分支时，完成时间取依赖图上关键路径的累积延迟，还需考虑同步与重试。各量单位为秒或明确的事件周期，不能与语义距离直接比较。平均延迟不足以保证截止期，实时任务应检查分位数、最坏界或失效概率；无法获得上界时，应诚实保留可用性不确定性。

延迟也会改变控制性质。考虑标量迭代 $x_{k+1}=x_k-a x_{k-d}$，$a>0$ 是归一化增益，$d$ 是整数延迟。无延迟时系数为 $1-a$，所以 $0<a<2$ 支持收缩；延迟为一步时特征多项式变成 $z^2-z+a$，例如 $a=1.5$ 的共轭根乘积为 $1.5$，两根模均为 $\sqrt{1.5}>1$，于是更新失稳。这个反例说明连续或无延迟的稳定参数不能直接用于远程工具、异步记忆和分布式执行。调度器应根据时间戳识别陈旧状态，必要时降低增益、等待确认或重规划。

记忆保留必须区分运行激活、内容存储和检索可达性。某个活动变量可以按 $\dot a=-\gamma a$ 衰减，其中 $\gamma$ 的单位是时间倒数，解为 $a(t)=a(0)e^{-\gamma t}$；当 $\gamma>0$，活动半衰期为 $\ln2/\gamma$。这只描述该变量的衰减，不说明文件、参数或情景事实同时被删除。磁盘记录可能长期保留却因索引失效不可召回，工作界面可能迅速释放却已将事实写入外部记忆。零衰减也不自动导致发散，是否累积失控还取决于输入、饱和、反馈和容量。

因此遗忘在工程上至少包含删除、降权、摘要、归档、失效标记和权限撤回。它们修改不同状态，拥有不同代价与回滚可能。AgentOS 的 $E$ 可以保留事件而调整召回权重，$\Theta$ 可以维持稳定规则但撤销失效依赖，$h(t)$ 可以释放当前对象而不损失长期记录。CCM 监测有限工作界面，Offloading 把程序、工具与资料外化到 $M_e$，Boundary 管理召回量和授权；三者不能靠一个介质阻尼系数全部解释。

探索尺度同样要从硬件噪声中分开。设有限候选评分为 $r_i$，采样分布 $p_i=\exp(r_i/\tau)/\sum_j\exp(r_j/\tau)$，$\tau>0$ 与评分使用相同尺度。尺度增加可以使分布更平，但若候选生成器从未包含正确动作，提高尺度也无济于事；若约束过滤失败，探索还会增加风险。芯片温度、随机数生成和软件采样参数是不同变量，只有经过具体测量和建模才能建立影响关系。TCE 的工程控制参数因此需要按分布和增益定义，不从“温度”这个名称推得物理退火结论。

能耗评估应在实际执行层记录。任务总焦耳可由功率随时间积分得到，调用费用和时延则各用独立计量；归一化加权可用于允许权衡的优化，但原始量应保留以便审计。生物与硅基系统不能仅由传导速度排列智能高低，因为并行性、缓存、感知、算法、学习和环境支持都会改变整体能力。硬件优势只支持特定接口和负载上的性能结论，不证明拥有更完整的语义结构或更高等级体验。

实现验证的设计因而围绕故障与边界展开。我们改变网络延迟、输入速率、内存配额和召回可用性，观察系统是否按模型降低更新频率、延迟非关键工作、请求资源或安全停止。若它在超预算时继续虚报成功，错误属于边界执行；若保留记录却无法检索，错误属于索引与读出；若采样尺度调整破坏授权，错误属于控制与规范耦合。这样的诊断比把所有现象称为“介质性格”更能指导修复，也更接近 HSF-HD 为工程提供一般状态语言的目标。

\section{状态空间几何：绑定、可辨识度量与操作次序}
\label{sec:cognitive_space_ontology}

本章最后说明状态空间的几何应怎样获得，而不是先宣布某种空间是认知的唯一物理本体。一个系统可以同时具有离散对象图、连续运行向量、概率信念和按时间组织的事件记录。整体空间可以是不同结构版本上的不交并；每个版本拥有自己的可行属性集合和读出接口。只有在固定版本内满足连续性与可微性时，才可使用梯度和局部度量。这样的混合状态组织既保留历史和结构事件，也避免把一切变化强行写成复数波函数。

张量绑定是其中一种表示设计。设角色向量 $u_i\in\mathbb R^r$，属性向量 $v_i\in\mathbb R^d$，可用 $T=\sum_i u_i v_i^\top\in\mathbb R^{r\times d}$ 保存角色—属性关联。若角色矩阵 $U=[u_1,\ldots,u_N]$ 列满秩，已知 $U$ 时可以通过左逆恢复相应属性矩阵；若角色不独立或身份未保存，多个绑定可能得到同一 $T$。张量积因此提供组合编码能力，却不证明量子纠缠，也不保证所有语义可无损恢复。对象身份、缺失值、误差和版本仍需独立记录。

局部结构网络与全局分布表示可以共同参与绑定。图 $G$ 保存实体、事件与依赖，向量 $z$ 保存语境与相似性，接口把二者的读出关联到同一任务对象。更新 $z$ 后不必自动修改 $G$ 的事实边；修改 $G$ 后也不保证向量空间同步保留全部差异。我们需要版本校验、冲突检测和有效期，防止召回向量绑定到已经删除的实体。what、where、when 可以作为属性、关系和时间来源的接口划分，却不把两种信息结构化约为三个固定神经通道。

信息几何需要额外概率假设。设 $p(o\mid\theta)$ 是参数 $\theta\in\mathbb R^q$ 的正则概率模型，支持集固定，导数与积分可以交换，评分 $s(o,\theta)=\nabla_\theta\log p(o\mid\theta)$。我们定义 Fisher 矩阵
\begin{equation}
 I(\theta)=\mathbb E_\theta[s(o,\theta)s(o,\theta)^\top].
\end{equation}
任意向量 $v$ 满足 $v^\top Iv=\mathbb E[(v^\top s)^2]\ge0$，所以它总是半正定；只有局部可辨识等额外条件才支持正定度量。两种参数若导致同一分布，评分在某方向为零，该方向不能被随意当作具有正长度的逻辑距离。

在上述正则条件及足够光滑性下，固定 $\theta$ 对 $\theta+\delta$ 的 KL 散度进行 Taylor 展开。归一化给出评分期望为零，一次项消失；二阶项由评分外积的期望形成，于是
\begin{equation}
 D_{\mathrm{KL}}\bigl(p_\theta\|p_{\theta+\delta}\bigr)
 =\frac12\delta^\top I(\theta)\delta+o(\|\delta\|^2).
\end{equation}
这说明局部度量来自可区分的概率分布变化，不说明它自动等于世界图距离、任务成本或社会价值。支持集改变、参数不可辨识或模型失配时，需要重新处理展开。自然梯度若使用 $I^{-1}$，还要保证可逆；阻尼或伪逆是算法选择，其含义不能被隐藏。

操作次序也可以产生非对易现象而无需量子公设。考虑一维状态上的映射 $A(x)=x+1$、$B(x)=2x$，则 $B(A(x))=2x+2$，而 $A(B(x))=2x+1$，次序显然影响结果。语义操作中的绑定、归档、权限撤回和结构编辑同样可能不交换，但必须按具体类型定义组合。有限维方阵也不能满足 $[Z,P]=i\hbar I$ 且 $\hbar\ne0$ 的正则关系，因为任何矩阵交换子的迹为零，右侧迹却为 $i\hbar n$。因此不能从次序敏感直接推出一种普遍语义不确定性原理。

有限容量带来的精度与覆盖权衡，可以用前节任务失真与预算表达；它不需要借用没有适用空间的对易式。如果确实采用无限维微分算子，还要声明定义域、边界条件和可观测量，相关数学结果只对那个专门模型成立。类似地，历史依赖可以通过事件序列、迟滞状态或慢变参数建模，不必称为有重量的几何体。保留概念之间的区别，才允许不同数学工具共同服务于状态动力学。

由此，我们获得本章的基本结论：智能状态模型由类型、身份、目标、结构、更新与边界共同构成；具体算子在声明条件内支持可证明性质；实现和现实任务通过独立证据验证。系统论、协同论、耗散理论、复杂系统理论与认知科学提供问题来源和约束，HSF-HD 构造一般模型族，AgentOS 将其细化为记忆与调度工程。后续章节会展开单纯复形、状态传播和多尺度适应，但每次扩展都必须继续遵守同一要求：先定义对象，再建立机制，最后限定结论和失效边界。
